\section{Conclusion and Future Work}
\label{sec:conclusion}

This work predicts, at the end of the previous day, whether a New York departure will arrive at
least 15 minutes late. It uses 59 leakage-safe features built from the published schedule and
past-only outcomes, Logistic Regression and XGBoost, and a strictly chronological protocol in
which each period is used for one decision. The test period was sealed before any analysis and
evaluated once. On it, the selected L2 Logistic Regression achieves PR-AUC 0.398 (delay rate 0.258)
and ROC-AUC 0.664. Sigmoid calibration reduces the expected calibration error from 0.204 to 0.077.
In Phase~2, a simple MLP, a GRU over daily delay sequences, an embedding MLP and a neural--XGBoost
hybrid, trained under the same protocol with five seeds each, did not rank flights better than
Logistic Regression beyond bootstrap noise. The GRU beat a control without recurrence, showing that
the order of recent daily delays carries signal, and scheduled congestion was the only feature group
whose removal reliably hurt. The GRU, frozen before the test and scored once, reached PR-AUC 0.399
and ROC-AUC 0.664, matching but not improving on Logistic Regression. Its September calibration
under-predicted on the test period (ECE 0.094), and its October threshold reached 57.4\% precision
there, against the 60\% target.

\textbf{Limitations.}
\begin{itemize}
  \item The provided dataset has no weather or same-day operational information, which bounds
  how confident any flight-level prediction can be.
  \item The validation months were calmer than training and test, so the calibration fitted on
  September under-predicts on the test period.
  \item The 60\% precision target is an assumption, not a measured cost trade-off, and was not
  reachable. No operating threshold or rescheduling rule could therefore be evaluated.
  \item History features assume flight outcomes are available within 30 minutes of arrival.
  \item One year of departures from three airports limits generalisation to other years and
  airports.
\end{itemize}

\textbf{Future work.} Adding weather forecasts available at $D{-}1$ is the most promising
extension. Recalibrating on recent data, such as a rolling recalibration window, would address the
seasonal shift in delay rates. With real cost estimates for false alarms and missed delays, the
operating point and the offline rescheduling rule could be set by expected cost. Since deeper models
did not add ranking quality on the same inputs, new information, not new architectures, is the
priority.
